\chapter*{Abbreviations and Acronyms}
\addcontentsline{toc}{chapter}{Abbreviations and Acronyms}
\markboth{Abbreviations and Acronyms}{Abbreviations and Acronyms}

\begin{tabularx}{\textwidth}{lX}
\toprule
\textbf{Term} & \textbf{Meaning} \\
\midrule
API   & Application Programming Interface \\
CRN   & Common Random Numbers \\
CSV   & Comma-Separated Values \\
DES   & Discrete-Event Simulation \\
DQN   & Deep Q-Network \\
FIFO  & First-In, First-Out \\
FTE   & Full-Time Equivalent \\
KPI   & Key Performance Indicator \\
MDP   & Markov Decision Process \\
RL    & Reinforcement Learning \\
RL-3  & The three-stage reinforcement learning sequencing agent evaluated in this thesis \\
SLA   & Service-Level Agreement \\
WIP   & Work In Progress \\
\bottomrule
\end{tabularx}

\vspace{1.0cm}
\noindent\textbf{Notation for workforce configurations.}
A workforce configuration is written as \rgm{sPKD}, where \(P\), \(K\) and \(D\) are the number
of monthly full-time equivalents allocated to picking, packing and dispatch respectively. For
example \rgm{s432} denotes four picking, three packing and two dispatch FTE. When any stage
requires ten or more FTE the unambiguous underscore form is used instead, for example
\rgm{s22\_11\_5} for 22 picking, 11 packing and 5 dispatch FTE.
